%! Radicals and Rational Exponents — Unit 1, Lesson 1.2 — Guided Notes (Answer Key)
\documentclass[10pt]{article}
\usepackage{atda-article}
\usepackage{atda-key}   % pulls in -boxes; do NOT also load -boxes
\newcommand{\TallMath}[1]{$\displaystyle #1\rule[-1.4em]{0pt}{3.2em}$}
\setlength{\workrowsep}{6pt}

\begin{document}

\pageheader{Unit 1, Lesson 1.2}{Guided Notes --- Answer Key}
% NAMESTRIP: no \namedateperiod here — mirrors the blank. Teacher-only prose goes
% in the lesson plan as \begin{teachernote}[Guided Notes], never in a key.

\vspace{0.15in}

\begin{objectivebox}
\small
By the end of this lesson, I will be able to\dots
\begin{enumerate}[leftmargin=1.5em, itemsep=2pt, topsep=3pt]
  \item simplify a radical with a \textbf{numeric or algebraic radicand} by pulling out every
        perfect $n$th-power factor;
  \item add, subtract, multiply, and divide radical expressions, including \textbf{rationalizing}
        a denominator, and simplify the result;
  \item convert between a radical and a \textbf{rational exponent} in both directions, and
        evaluate expressions like $27^{2/3}$ without a calculator;
  \item use a square-root model from a real situation --- speed from a skid mark ---
        to answer a question and say what the answer means.
\end{enumerate}
\end{objectivebox}

\vspace{0.10in}

\boxguard
\begin{vocabbox}
\small
Fill in each term as we build it together.
\vocabans{Radical, radicand, and index}{In $\sqrt[n]{a}$, the whole symbol is the radical, $a$ is
the radicand (what is inside), and $n$ is the index (which root). A missing index means $2$.}
\vocabans{Principal square root}{The non-negative square root. $\sqrt{49}$ means $7$, not $-7$ ---
the negative one is written $-\sqrt{49}$.}
\vocabans{Simplified radical form}{No perfect $n$th-power factor remains in the radicand, and no
radical remains in a denominator.}
\vocabans{Rational exponent}{A fractional exponent: $a^{m/n}=\sqrt[n]{a^m}$. The denominator is
the index of the root; the numerator is the power.}
\vocabans{Rationalizing the denominator}{Multiplying a fraction by a form of $1$ that clears the
radical out of the denominator, e.g. $\dfrac{6}{\sqrt3}\cdot\dfrac{\sqrt3}{\sqrt3}=2\sqrt3$.}
\end{vocabbox}

\vspace{0.10in}

\boxguard[12]
\begin{hookbox}
\small
\textbf{Reading a skid mark.} Crash investigators estimate how fast a car was going from the
length of its skid marks. On a road with drag factor $f$, a skid of $d$ feet means the car was
travelling about
\[ s = \sqrt{30fd} \quad \text{miles per hour.} \]
Two cars skid to a stop on the same dry asphalt ($f=0.8$). Car A leaves a $60$-foot skid;
Car B leaves a $240$-foot skid --- \textbf{four times as long}.

Was Car B going four times as fast? Decide first, then check it.

\ansline{Car A: $\sqrt{1440}=12\sqrt{10}\approx 37.9$ mph. Car B: $\sqrt{5760}=24\sqrt{10}\approx 75.9$ mph.}\\
\ansline{No --- four times the skid gives only \emph{twice} the speed, because $\sqrt4=2$.}\\
\end{hookbox}

\vspace{0.10in}

\boxguard[24]
\begin{notesbox}{1. Simplifying a Radical --- Pull Out the Perfect Powers}
\small
$\sqrt[n]{a}$ asks: what number, raised to the $n$th power, gives $a$? The $n$ is the
\textbf{index}, and $a$ is the \textbf{radicand}. Two properties do all the work:
\[ \sqrt[n]{ab}=\sqrt[n]{a}\cdot\sqrt[n]{b} \qquad\text{and}\qquad
   \sqrt[n]{\dfrac{a}{b}}=\dfrac{\sqrt[n]{a}}{\sqrt[n]{b}} \quad (b \ne 0). \]
So: find the largest perfect $n$th-power factor inside, take its root, and leave the rest.

\renewcommand{\arraystretch}{1.9}
\begin{tabularx}{\linewidth}{@{}p{4.4cm} p{4.6cm} X@{}}
\rowcolor{mist}
\textbf{Radical} & \textbf{Split off the perfect power} & \textbf{Simplified} \\
\hline
$\sqrt{50}$          & $\sqrt{25 \cdot 2}$        & \ans{$5\sqrt{2}$} \\
$\sqrt[3]{54}$       & $\sqrt[3]{27 \cdot 2}$     & \ans{$3\sqrt[3]{2}$} \\
$\sqrt{72x^5}$       & $\sqrt{36x^4 \cdot 2x}$    & \ans{$6x^2\sqrt{2x}$} \\
$\sqrt[3]{40y^7}$    & $\sqrt[3]{8y^6 \cdot 5y}$  & \ans{$2y^2\sqrt[3]{5y}$} \\
$\sqrt{\dfrac{49}{16}}$ & $\dfrac{\sqrt{49}}{\sqrt{16}}$ & \ans{$\dfrac{7}{4}$} \\
\end{tabularx}

\vspace{6pt}
\textbf{Careful:} a radical is only simplified when \emph{no} perfect $n$th-power factor is left
inside. Splitting $\sqrt{72}$ as $\sqrt{4}\cdot\sqrt{18}$ is true but not \ans{simplified} --- there
is still a perfect square hiding in the $18$.

\textbf{Back to the hook} --- Car A's skid, exactly and then as a decimal:
\begin{work}
  s &= \sqrt{30(0.8)(60)} \\
    &= \sqrt{1440} \\
    &= \sqrt{144 \cdot 10} \\
    &= 12\sqrt{10} \approx 37.9 \text{ mph}
\end{work}
\end{notesbox}

\vspace{0.10in}

\boxguard[22]
\begin{notesbox}{2. Rational Exponents --- Nobody Chose These Either}
\small
In Lesson 1.1 the zero and negative exponents were \emph{forced} on us by the quotient rule.
Fractional exponents are forced the same way, by the product rule:
\[ 9^{1/2}\cdot 9^{1/2} = 9^{\frac12+\frac12} = 9^1 = 9, \]
so $9^{1/2}$ has to be the number that squares to $9$. That is exactly what $\sqrt{9}$ means.

\begin{itemize}[leftmargin=1.3em, itemsep=3pt, topsep=3pt]
  \item $a^{1/n} = \sqrt[n]{a}$ \quad --- the denominator of the exponent is the
        \ans{index} of the radical.
  \item $a^{m/n} = \sqrt[n]{a^m} = \left(\sqrt[n]{a}\right)^m$ \quad --- the numerator is the
        power; take the root first, since the numbers stay smaller.
\end{itemize}

\renewcommand{\arraystretch}{1.8}
\begin{tabularx}{\linewidth}{@{}p{4.2cm} X | p{3.6cm} X@{}}
\rowcolor{mist}
\textbf{Radical form} & \textbf{Rational exponent} & \textbf{Evaluate} & \textbf{Value} \\
\hline
$\sqrt[3]{x^5}$          & \ans{$x^{5/3}$} & $27^{2/3}$   & \ans{$9$} \\
$\sqrt[4]{y^3}$          & \ans{$y^{3/4}$} & $16^{-1/2}$  & \ans{$\dfrac14$} \\
$\sqrt{t}\cdot\sqrt[3]{t}$ & \ans{$t^{5/6}$} & $81^{3/4}$ & \ans{$27$} \\
\end{tabularx}

\vspace{6pt}
\textbf{Every exponent rule from Lesson 1.1 still holds} --- the exponents just became fractions:
\begin{work}
  \left(32x^{10}\right)^{3/5} &= 32^{3/5}\left(x^{10}\right)^{3/5} \\
                              &= \left(\sqrt[5]{32}\right)^3 x^{6} \\
                              &= 8x^6
\end{work}
\end{notesbox}

\vspace{0.10in}

\boxguard[20]
\begin{notesbox}{3. Adding, Subtracting, Multiplying, and Dividing Radicals}
\small
\textbf{Adding and subtracting: like radicals only.} Two radicals combine when they have the same
index \emph{and} the same \ans{radicand} --- exactly the like-terms rule from Lesson 1.1, with
$\sqrt{3}$ playing the part of $x$. Simplify first; radicals that look unlike often are not.
\begin{work}
  3\sqrt{12}+\sqrt{75} &= 3\left(2\sqrt{3}\right)+5\sqrt{3} \\
                       &= 6\sqrt{3}+5\sqrt{3} \\
                       &= 11\sqrt{3}
\end{work}

\textbf{Multiplying: every term times every term}, then simplify what lands inside.
\begin{work}
  \left(2\sqrt{5}\right)\left(3\sqrt{10}\right) &= 6\sqrt{50} \\
                                                &= 6 \cdot 5\sqrt{2} \\
                                                &= 30\sqrt{2}
\end{work}
\begin{work}
  \left(\sqrt{7}+3\right)\left(\sqrt{7}-2\right) &= 7-2\sqrt{7}+3\sqrt{7}-6 \\
                                                 &= 1+\sqrt{7}
\end{work}

\textbf{Dividing: rationalize the denominator.} A simplified answer carries no radical downstairs.
Multiply top and bottom by whatever makes the denominator a perfect power --- you are multiplying
by $1$, so the value never changes.
\begin{work}
  \dfrac{6}{\sqrt{3}} &= \dfrac{6}{\sqrt{3}}\cdot\dfrac{\sqrt{3}}{\sqrt{3}} \\
                      &= \dfrac{6\sqrt{3}}{3} \\
                      &= 2\sqrt{3}
\end{work}
\end{notesbox}

\vspace{0.10in}

\boxguard[20]
\begin{practicebox}
\small
\textbf{The accident report.} A car skids to a stop on \emph{wet} asphalt, where the drag factor
drops to $f=0.5$. The investigator measures the skid at $d=120$ feet. The posted limit on that
road is $35$ mph.

\begin{enumerate}[leftmargin=1.6em, itemsep=6pt, topsep=4pt]
  \item Use $s=\sqrt{30fd}$ to find the speed. Give the exact simplified radical, then round to
        the nearest tenth.
\begin{work}
  s &= \sqrt{30(0.5)(120)} \\
    &= \sqrt{1800} \\
    &= \sqrt{900 \cdot 2} \\
    &= 30\sqrt{2} \approx 42.4 \text{ mph}
\end{work}

  \item Was the driver speeding? Answer with a number and a sentence.

\ansline{Yes --- about $42.4$ mph in a $35$ mph zone, roughly $7$ mph over the limit.}

  \item The investigator wants to know how long a skid would mean the car was going
        \emph{twice} this fast. Explain, using the radical, why the answer is $480$ feet and not
        $240$ feet.
\begin{work}
  \sqrt{30(0.5)(480)} &= \sqrt{7200} \\
                      &= 60\sqrt{2} \approx 84.9 \text{ mph}
\end{work}
\ansline{Speed depends on $\sqrt d$, so doubling the speed needs $2^2=4$ times the skid length.}\\
\ansline{Doubling $d$ to $240$ ft only multiplies the speed by $\sqrt2 \approx 1.4$.}\\
\end{enumerate}
\end{practicebox}

\end{document}
